\maketitle
\head{本版编排与收录说明}
\noindent\textbf{这是一份按内容顺序整理的中文工作译稿，不是全部已上传材料的完结全译本。}第I、II章及第III章的两个已出版分册正文已经连续收入；第0章当前收入§§1--19及§20的20.1--20.5，止于（20.5.15）。20.6--20.7、§§21--23、第IV章主体及第V卷预备笔记余文尚未连续译入。未译内容没有用原文扫描页、提要或空白标题代替译文。

\section*{正文次序}
正文按原著内容编号组织：导言；第0章预备知识；第I章概形的语言；第II章若干类态射的初等整体研究；第III章凝聚层的上同调研究；后续待译范围说明；独立编排的第V卷预备笔记已译部分；最后为各原分册的书末资料。第III章第一部分的§§1--5与第二部分的§§6--7直接衔接。

第0章原来分载于第一卷、第三卷第一部分、第四卷第一部分。本版将已译补篇归并在同一章，但保留 $0_{\mathrm I}$、$0_{\mathrm{III}}$、$0_{\mathrm{IV}}$ 等出处记号；它们标示原出版分册，并非重复的三章。这一安排依内容编号，不声称每节的全部前提均已在前文出现；原著跨章引用仍按原号保留。

\section*{实际收录范围}
\begin{longtable}{@{}p{.25\linewidth}p{.67\linewidth}@{}}
\toprule 材料 & 已收入的连续译文及状态\\\midrule\endhead
第一卷（1960） & 导言；第0章§§1--7；第I章§§1--10，止于（10.15.7）；原书书末资料。承接既有V2校订源码。\\
第二卷（1961） & 第II章§§1--8，止于（8.14.14）；参考文献、索引、目录、勘误清单一。书末资料沿用此前依据底本恢复的译排来源。\\
第三卷第一部分（1961） & 第0章§§8--13；第III章§§1--5，止于（5.5.2）；原书书末资料。\\
第三卷第二部分（1963） & 第III章§§6--7，止于（7.9.14），连同证明、记号索引、术语索引、原目录及勘误清单二。\\
第四卷第一部分（1964） & 第0章补篇章首、§§14--19及20.1--20.5，止于（20.5.15）。（19.7.1.5）较弱原版的拓扑范围仍待核；原文与已核定形式的差别保存在正文脚注及报告。\\
第V卷英文预备笔记 & \texttt{EGAV1-2.ps}六页可读材料；\texttt{EGAVaiello.pdf}卷首说明及“超平面截面与锥投影”§§1--6。余文不计入已译范围。\\
\bottomrule
\end{longtable}

\section*{正文、脚注与核查报告}
新增正文采用已经核实的正确陈述、公式和证明；原文错误写法、修改理由及依据放在脚注。作者正式勘误与译者补证分开，较长反例与验证收入独立核查报告。未确认的疑点不擅自定为错误；原著本来省略的证明、预备笔记中未写完或无法辨认的内容，不凭常识补造成“原文”。本轮新增注号以I9开头，旧注号保留。

这一原则适用于本轮新增20.1--20.5，并承接第七、八版明确列出的回溯校订。\textbf{它不表示此前全部旧译文中的历史核查意见都已经逐项改入正文。}其余既有数学正文保留，校订状态须结合相应脚注与历史报告阅读。对作者已经撤销的（19.3.12），正文继续只保留编号与撤销状态；完整历史陈述及证明在核查报告中存档。

第V卷上传载体属于预备笔记的翻译与编辑，英译编者为Piotr Blass与Joseph Blass。重组建议、未完成句子和待补证明不按正式刊行的第I--IV章同等状态标示。与英文合并稿重合的PS载体不重复计作新译文。

\section*{第九版新增范围与核查边界}
本版从§20“导子与微分”章首开始，连续译入20.1--20.5：导子与扩张、函子正合列、连续导子、主部和微分模、微分模的基本函子性质。共49个原著编号条目、57个原编号公式或交换图标签。对应第四卷第一部分分册印刷第116--136页的相关部分（上传PDF第115--135页）；下一接续点为20.6“非完美性模与特征同态”及（20.6.1）。

（20.4.13，ii）原同构须改为沿增广与基环张量后的同构；（20.5.4.2）一般自然性图的源对象须校正；（20.3.2）的连续推广需要明确拓扑下降条件。本版分别给出反例、正确陈述或充分条件，并与作者关于标量环、外微分符号及整闭性假设的正式修订分开。它们不是同一性质的“错误定理”。

本轮仍不解决旧版H8-P01等已标出的待核事项，也不宣称对全书旧译文重新完成逐句审校。

\section*{版式与使用}
采用 \texttt{\string\documentclass[UTF8,a4paper]\{ctexbook\}}，由文档类决定默认中文字体、字号、版心、标题及双面分页。不另设页边距，不用自定义页眉。原章号、条目号与公式号保留；默认双面排版所需的空白偶数页不删除。前置部分用罗马数字页码，正文及书末资料连续编号。总目录指向本合订本页面；各原著目录的页码仍指原版。

PDF、可独立编译的单文件TeX和分文件工程一并提供。实际编译与检查结果见工程包中的记录。编译、条目覆盖及渲染一致性不代替出版级逐句审校或全书证明的独立验证。
\tableofcontents
